\chapter{Building interactive charts and infographics}\label{ch:interactive}

The board wants \enquote{something like the newspaper does,} an interactive map they can hover, a chart that animates, and they want it for the quarterly meeting. Static figures are right for a printed report, but a client who will explore the data, or an audience you want to give a tool they keep, calls for something they can touch. The risk is that interactivity leaves the replicable pipeline and becomes a designer's one-off; in this chapter we keep it inside Stata.

In this chapter we build small panels for comparing trends, an interactive state map, a searchable table, and an animated chart. We then test whether the intended readers can use the pages and document how to refresh them. You will learn to choose a form that fits the reader's task and preserve the source data, code, and dated output needed to reproduce a release.\marginnote{Some chart pages fetch a plotting library from an outside server when opened. Test the page on the client's network before delivery, because a blocked request can prevent the chart from appearing.}

\section{The visualization suite: one toolkit, seven tools}\label{sec:suite}
Where does a deliverable like the board's interactive map actually come from? In this book, from a short chain of commands: across this chapter and the two that follow we draw on seven of our own Stata packages, built to work together, each a single command that takes a Stata dataset one step closer to something a client can open in a browser. The packages cover different stages of the work. Two of them bring data \emph{in}, four \emph{render} it as a chart or dashboard, and one \emph{wraps} the results into a finished report or web portal. Once you know where each tool sits in that pipeline, you assemble the next deliverable from parts instead of hand-building it every time.

Two terms will come up on every page of this chapter, so pin them down now. An \emph{interactive chart}\index{interactive chart} is a picture the reader can act on, hover a state to read its value, type a name into a search box, drag a slider, rather than only look at. \emph{Self-contained HTML}\index{self-contained HTML} means the whole thing, the data and the code that draws it, is stored inside one \texttt{.html} file you can email, so \enquote{it renders in the browser} simply means the reader double-clicks that file and the chart draws itself on their screen, no software to install.\marginnote{HTML is the plain-text format web pages are written in; a browser is any program that opens it (Chrome, Safari, Edge). \enquote{Renders} is just the browser's word for \enquote{draws.} If you can open a web page, you can open everything this suite produces.}

We sort the seven tools with two questions. \emph{Online versus offline}: does the finished page need an internet connection to draw itself, or does it contain everything it needs? \emph{Chart versus container}: is the tool drawing a single graphic, or is it a shell that combines many graphics and prose? Figure~\ref{fig:suite-pipeline} lays out the pipeline, and Table~\ref{tab:suite-matrix} places each tool in the grid those two axes make.

\begin{figure*}[htbp]
\centering
\begin{tikzpicture}[node distance=0.5cm and 0.75cm]
  \node[stage] (ing) {Ingest};
  \node[stage, right=of ing] (ren) {Render};
  \node[stage, right=of ren] (wrap) {Wrap};
  \draw[flow] (ing) -- (ren);
  \draw[flow] (ren) -- (wrap);
  \node[substep, below=0.35cm of ing, text width=2.0cm]
    {\texttt{webapi} (Ch.\,3)\\\texttt{googlesheets} (Ch.\,11)};
  \node[substep, below=0.35cm of ren, text width=3.0cm]
    {\texttt{googlechart} (online)\\\texttt{sparkta2} (offline)\\\texttt{statashiny} (Ch.\,12)\\\texttt{dashboardbuilder} (Ch.\,12)};
  \node[substep, below=0.35cm of wrap, text width=2.0cm]
    {\texttt{webdoc2} (Ch.\,12)\\one report or portal};
\end{tikzpicture}
\caption[The visualization suite pipeline]{The suite as one pipeline. Data comes in through \texttt{webapi} (JSON APIs, Chapter~\ref{ch:apis}) or \texttt{googlesheets} (live Google Sheets, Chapter~\ref{ch:sheets}); it is rendered as a chart by \texttt{googlechart} (online) or \texttt{sparkta2} (offline), or as an interactive dashboard by \texttt{statashiny} or \texttt{dashboardbuilder}; and \texttt{webdoc2} wraps the results into one shippable report or GitHub-Pages portal (all Chapter~\ref{ch:portals}). We cover the render column's two chart tools in this chapter.}
\label{fig:suite-pipeline}
\end{figure*}

\subsection{Pick the tool: online or offline, chart or container}\label{sec:suite-choose}
Faced with a tool choice, ask the grid's two questions and let the answers pick for you. Does the page have to work with the internet off? Then you are in the offline row: \texttt{sparkta2}\index{sparkta2@\texttt{sparkta2}} for a chart, \texttt{statashiny}\index{statashiny@\texttt{statashiny}}, \texttt{dashboardbuilder}\index{dashboardbuilder@\texttt{dashboardbuilder}}, or \texttt{webdoc2}\index{webdoc2@\texttt{webdoc2}} for a container. Are you drawing one graphic or combining many? That is the chart-versus-container column. Those two questions narrow the choice to a single tool everywhere except the offline-container cell, which has three. \texttt{webdoc2} is the wrapper you reach for when the deliverable is a whole report; between the two dashboards, pick by the reader's verb: \texttt{statashiny} when readers explore rows with a search box or slider, \texttt{dashboardbuilder} when readers scan headline tiles against a benchmark (Chapter~\ref{ch:portals} works all three). Name the venue and the deliverable, look up the cell where those two answers meet, and where that cell lists more than one tool, choose by what the reader will actually do.

Spend one more minute with the grid before any command runs. Ahead of the build, write the deliverable in one sentence that answers both axis questions, something like \enquote{a county wage map for the agency website, refreshed quarterly, firewall status unconfirmed}: the venue clause fixes the row, the refresh clause is where you see that the page is a standing commitment (Section~\ref{ch10-sec:maintenance}), and the unconfirmed firewall is the reason to test with the network off before promising the online tool. Write that sentence where the team can see it, and nobody has to relitigate the tool choice at every revision.

Two tools in the grid are the pair readers most often confuse, because both have Google in the name and they behave nothing alike. \texttt{googlechart}, the subject of this chapter, and \texttt{googlesheets}, the subject of Chapter~\ref{ch:sheets}, differ on three things worth knowing before you commit to either. They differ first on \emph{credentials}: \texttt{googlesheets} acts as you inside your own Google account, so it needs the one-time browser sign-in of Appendix~\ref{app:setup}, while \texttt{googlechart} writes a local file and needs no Google account at all. If web charts are all you want, you never open the Google Cloud console. They differ second on \emph{who reads the result}: a \texttt{googlesheets} write goes straight into a document your collaborators are already editing, so you and they work from one copy, whereas a \texttt{googlechart} page is a file you publish, which anyone with the link can read, filter, and export from but nobody can change. Collaborators need a workbook they can keep editing; a public audience needs a figure that will not shift under them. Those two differences explain the third, why they fail differently, which Section~\ref{sec:gc-how} works out for \texttt{googlechart} and Chapter~\ref{ch:sheets} for \texttt{googlesheets}. When a project needs both halves, you run the two commands in the same do-file.

\begin{table*}[htbp]
\centering
\caption[The visualization suite at a glance]{The seven-tool suite placed on the two axes that decide between them: online versus offline (does the finished page need the internet to draw?) and chart versus container (one graphic, or a shell that combines many?). \texttt{webapi} sits in the ingest column and is shown for completeness; the render and wrap tools are the subject of this chapter and the next two. Read down the \enquote{Best for} column to pick a tool by the job in front of you.}
\label{tab:suite-matrix}
{\footnotesize
\setlength{\tabcolsep}{4pt}
\begin{tabular}{lllll}
\hline
Tool & Layer & Renders & Net at view & Best for \\
\hline
\texttt{webapi}      & Ingest    & (data in)      & n/a  & JSON APIs feeding the suite (Ch.\,3) \\
\texttt{googlesheets}& Ingest    & (data I/O)     & live & Live Google Sheets as source or surface (Ch.\,11) \\
\texttt{googlechart} & Chart     & Online (Google)& yes  & Web embeds, filter dropdowns, US-state geo, bubble-with-Play \\
\texttt{sparkta2}    & Chart     & Offline (D3)   & no   & Air-gapped briefs; TX county / district maps \\
\texttt{statashiny}  & Container & Offline (JS)   & no   & Single-file dashboards: sliders, search, cards (Ch.\,12) \\
\texttt{dashboardbuilder} & Container & Offline (SVG) & no & KPI tiles, tabs, benchmark selector (Ch.\,12) \\
\texttt{webdoc2}     & Container & Offline (BS-5) & no   & The report/portal shell that embeds the above (Ch.\,12) \\
\hline
\end{tabular}
}
\end{table*}

In this chapter we build out the two chart tools in the render column, \texttt{googlechart}\index{googlechart@\texttt{googlechart}} online and \texttt{sparkta2} offline\marginnote{A chart is the atom of every deliverable the suite ships, which is why the render column's two chart tools come first.}. Chapter~\ref{ch:sheets} covers \texttt{googlesheets}\index{googlesheets@\texttt{googlesheets}}, the live-data connector, and Chapter~\ref{ch:portals} covers the three containers, \texttt{statashiny}, \texttt{dashboardbuilder}, and \texttt{webdoc2}, that wrap these charts into a single dashboard, report, or portal. Keep the grid in mind as we go: every command below is an instance of one cell in it.

\section{Sparklines: the word-sized trend}
When a portfolio has forty sites, forty separate trend charts is not communication, it is a stack no one reads. A sparkline\index{sparkline}, a word-sized line with no axes, solves this by fitting a whole trajectory into a table cell,\marginnote{\textcite{tufte2006beautiful} coined \enquote{sparkline} in \emph{Beautiful Evidence}, defining it as a \enquote{small, intense, word-sized graphic.} The point is that it sits inline with text, at the resolution of the surrounding type.} so a reader scans one column and sees every site's trend at once, the rising ones and the falling ones jumping out. The \texttt{sparkta2} engine generates these inline trends and places them in dashboard tables.

\begin{kaobox}[frametitle=Package Integration: \texttt{sparkta2}]
\textbf{Integrated command:} \texttt{sparkta2}. The suite's \emph{offline} chart tool (Table~\ref{tab:suite-matrix}): an engine built on D3, the JavaScript charting library many newsroom graphics teams build on, that generates high-density inline trend sparklines and sub-state maps, rendering fully offline.
\end{kaobox}

One honest caveat applies to this tool. As of this writing the \texttt{sparkta2} source is not yet published, only rendered galleries are public, so read this section as a description of what the engine renders, not an install you can run today.

\subsection{The spark wall you can build today}\label{sec:sparkwall}
You do not have to wait for \texttt{sparkta2} to get the sparkline's payoff, because native Stata already draws small multiples\index{small multiples}\marginnote{Small multiples: a wall of tiny panels that strip their axes to almost nothing so the shape is all that is left to read.}. A regional workforce board wants to know whether the wage recovery after 2020 looked the same across its labor markets, or whether some counties bounced and others stalled, because it targets its programs county by county. Six separate wage charts would be six page-turns; a spark wall puts all six side by side so the divergence is one glance.

We build the wall from the same no-key BLS files Chapter~\ref{ch:apis} introduced, assembling the county-quarter panel that chapter promised: one \texttt{import delimited} per county-quarter, looped over six Texas counties and twenty-four quarters. The Quarterly Census of Employment and Wages posts one CSV per county-year-quarter at a stable URL, and we keep the total private-sector line (ownership code~5, industry code~10) from each:
\par
{\footnotesize
\begin{verbatim}
local counties 48453 48201 48113 48029 48439 48141
tempfile master
local first 1
foreach fips of local counties {
  forvalues yr = 2019/2024 {
    forvalues q = 1/4 {
      local url ///
        "https://data.bls.gov/cew/data/api/`yr'/`q'/area/`fips'.csv"
      capture import delimited "`url'", ///
        clear varnames(1) stringcols(1 3)
      if _rc continue
      keep if own_code == 5 & industry_code == "10"
      keep area_fips year qtr avg_wkly_wage
      if `first' save `master', replace
      else       append using `master'
      save `master', replace
      local first 0
    }
  }
}
\end{verbatim}
}
Before you adapt this loop, notice two details. The \texttt{capture} before \texttt{import} and the \texttt{if \_rc continue} after it keep one missing quarter from aborting the whole loop: Stata skips that quarter and carries on, which matters because the newest quarters lag by several months and one county may not have posted yet.\marginnote{Robust loops over live URLs need a skip clause: \texttt{capture} swallows the error and \texttt{continue} moves to the next iteration, so one 404 out of 144 downloads costs you one panel, not the whole run. Check \texttt{\_N} after to see how many landed.} And we again read columns~1 and~3 as strings with \texttt{stringcols(1 3)}, because the FIPS code and the industry code are identifiers, not quantities, the same leading-zero habit that Chapter~\ref{ch:apis} argued prevents silent merge failures. When you do adapt it, pilot the loop first with one county and one year, confirm the import returns four rows with the wage column populated, and only then widen the locals to the full list.

Once stacked, we build a quarter index and let \texttt{by()} draw the wall. The trick that makes it a spark wall rather than six ordinary charts is stripping the axes: tiny labels, no gridlines, no axis titles, so the eye reads the shape and not the furniture.
\par
{\footnotesize
\begin{verbatim}
gen tq = yq(year, qtr)
format tq %tq
twoway line avg_wkly_wage tq, ///
    lcolor(navy) lwidth(medthin) ///
    by(county, legend(off) rows(2) ///
       note("BLS QCEW, no key")) ///
    ytitle("") xtitle("") ///
    ylabel(, nogrid labsize(vsmall)) ///
    xlabel(, nogrid labsize(vsmall))
graph export "$figures/ch10_sparkwall.png", ///
    replace width(2400)
\end{verbatim}
}
The panel is real, not simulated: 144 attempted downloads, one row kept per county-quarter. Before you trust any figure built from live URLs, check it against something you already know, and this wall passes that check. Every county climbs across the window, and the first-quarter spike then dip repeats in all six panels, the sawtooth that annual bonuses leave in quarterly wage data.\marginnote{The levels sit in the axis ranges: Travis, Dallas, and Harris, the tech-and-corporate markets, top out near \$2{,}000 a week, while El Paso never clears \$900 and Bexar sits in the low \$1{,}200s.} A board member scanning the wall takes in the whole portfolio at once, the recovery lifted every county on the same seasonal rhythm without closing the gap between them, and no animated chart would deliver that comparison faster.

\begin{figure*}[htbp]
\centering
\includegraphics[width=\linewidth]{ch10_sparkwall.png}
\caption[A spark wall of Texas county wages]{Private-sector average weekly wage for six Texas counties, 2019q1--2024q4, built end to end by \texttt{ch10\_smallmult.do} from public BLS QCEW area CSVs with no API key. Each panel is a stripped-down sparkline on its own free y-scale, so the shape of each county's trajectory reads even where the wage level differs; the command that made it is printed above. Add a FIPS code to the \texttt{counties} local and the wall grows a panel. Six trajectories read at once, so a reader scans the whole portfolio without turning a stack of pages.}
\label{fig:sparkwall}
\end{figure*}

The whole exhibit is one short do-file (\texttt{ch10\_smallmult.do}) that reruns from the same public URLs, so next quarter's data point arrives by editing the year range in one \texttt{forvalues} line. Because the density is generated rather than hand-placed, the wall refreshes when the data does instead of drifting out of date in a design file; because the counties are listed in one local, adding a county grows a panel; and because a board reads it in the dollars a week and the county names it already uses, nobody needs you in the room to interpret it.

Notice which lines of the do-file do the steering. Everything that will change next quarter is defined in exactly two places, the \texttt{counties} local and the \texttt{forvalues yr} range, and everything below them is machinery that never needs touching. That separation is worth building on purpose: hoist every value you expect to edit, the geography list, the year range, the output path, into a header block of locals at the top of the file, and the quarterly refresh becomes a one-line edit a successor can make without reading the loop.\marginnote{A quick test of whether you have parameterized enough: ask what next quarter's update requires; if the answer names any line below the header, move that value up.}

\subsection{What \texttt{sparkta2} adds: offline maps (source not yet released)}\label{sec:sparkta2}
The spark wall is native Stata; \texttt{sparkta2} is the suite tool that generalizes it, and it occupies the one grid cell \texttt{googlechart} cannot reach: an interactive chart that draws with the internet off, and a map below the state level. Where \texttt{googlechart}'s geographic type stops at US states (Section~\ref{sec:gc-geo}), \texttt{sparkta2} renders Texas \emph{county} and even \emph{school-district} choropleths, the sub-state geography an education or workforce evaluator actually reports in. It does this by bundling its drawing engine (D3, a JavaScript charting library) \emph{inside} the HTML file, so there is no live library to fetch and nothing to blank behind a firewall.\marginnote{\enquote{D3} is a widely used JavaScript library for data-driven graphics. The point for us is not the library but where it lives: \texttt{sparkta2} ships it inside the file, so the page is truly self-contained. \texttt{googlechart} instead links out to Google's copy at view time, which is the online-versus-offline split of Table~\ref{tab:suite-matrix}.}

Its planned type list reads like the offline complement to \texttt{googlechart}: multi-series lines, diverging bars, donuts, a bar-race animation, hexbin density maps, bivariate two-variable choropleths, and the sub-state maps that are its signature. The decision between the two siblings comes down to two questions, both drawn straight from Table~\ref{tab:suite-matrix}. Must the page render with no internet, for an air-gapped agency or an embargoed report? Choose \texttt{sparkta2}. Do you need geography finer than the state, a map of Texas counties or independent school districts? Choose \texttt{sparkta2}, because \texttt{googlechart}'s \texttt{geo} type cannot draw it. For everything else, native filter dropdowns, an animated bubble with a Play button, a searchable table, US-state maps, \texttt{googlechart} is the richer tool, and we spend the rest of this chapter on its worked examples. The caveat in the heading still applies: \texttt{sparkta2} is not yet published.

\section{Fourteen chart types from one command}
You hand a client a page they can explore rather than a figure they can only look at, and you build it without leaving the do-file. Four worked examples follow, a state choropleth, an animated bubble with a Play button, a searchable table, and a diverging Likert bar, and once you have built those four you will be able to reach the other ten chart types from the same grammar, so answering a client's next question means changing one option rather than starting a new build. The \texttt{googlechart} command turns a Stata dataset into one interactive HTML page, with fourteen chart types from one syntax: geographic choropleths, bubbles that animate across a time panel, searchable tables, diverging Likert bars, and the rest. It includes a brand theme and an optional download menu, so a client can export the underlying data or a PNG themselves, which turns a static deliverable into a small tool the audience can use.

Because the chart is written by a do-file, it stays inside the replicable pipeline: the interactive graphic rebuilds when the data updates, rather than living as a hand-made artifact in a design file that drifts out of date. That is the key difference from a one-off dashboard, and it is why infographic-quality output belongs in code.\marginnote{A useful test of any \enquote{interactive} deliverable: could you regenerate last quarter's version from scratch if you had to? If the answer is \enquote{only by re-clicking through a design tool,} it is not really in the pipeline yet.}

\begin{kaobox}[frametitle=Package Integration: \texttt{googlechart}]
\textbf{Integrated command:} \texttt{googlechart}. The suite's \emph{online} chart tool (Table~\ref{tab:suite-matrix}). From any dataset, writes one interactive HTML page (column, bar, line, area, combo, pie, donut, scatter, animated bubble, geo, timeline, searchable table, histogram, and diverging stacked bar), with filter dropdowns, a brand theme, and an optional data/PNG download menu. It writes the file with no key and no network; the browser fetches Google's drawing code only when the page is opened.\marginnote{Requirements: \texttt{googlechart} needs Stata~17 or later and calls no Python, so there is no separate environment to set up before you run it.} It is published on SSC, so one line installs it:
\par
{\footnotesize
\begin{verbatim}
ssc install googlechart
\end{verbatim}
}
\end{kaobox}

\subsection{How one command writes an interactive page}\label{sec:gc-how}
The \enquote{no key} promise and the \enquote{needs network at view time} caveat sound like a contradiction until you see what actually happens. When you run \texttt{googlechart}, it does not call any web service. It writes a single text file that contains three things: your data, embedded as JSON (the labeled-box text format of Chapter~\ref{ch:apis}); a small block of styling; and, at the bottom of the file, the short JavaScript that hands both to the chart library and asks it to draw. Nothing is fetched while Stata runs, so there is no key, no login, and no network call at build time. The one thing the file does \emph{not} contain is Google's chart-drawing library itself, which its license forbids redistributing, so the page links out to Google's copy. The reader's browser follows that link when they open the page, which is what \enquote{needs network at view time} means: the file builds offline but draws online.\marginnote{This is the online cell of Table~\ref{tab:suite-matrix} made concrete. The file contains your data and your logic; only the last library is Google's and is fetched on open. If the reader's network blocks \texttt{gstatic.com}, the page blanks, which is exactly the failure the offline sibling \texttt{sparkta2} avoids.}

Every example below follows one shape. Load a dataset, call \texttt{googlechart} with a \texttt{type()} and a few options, and it writes one \texttt{.html} file named by \texttt{export()}. The \texttt{noopen} option we use throughout just suppresses the auto-open so a do-file can build a dozen pages without popping a dozen browser tabs. These commands are runnable: the companion \texttt{ch10\_googlechart.do} builds all six pages in this chapter, offline and with no key, and you can open the results in any browser. The \texttt{DATA} and \texttt{OUT} locals in each listing are the same header-block habit as the spark wall's county list: the do-file defines both paths once at the top, so repointing all six builds at a new data release or a different output folder is one edit, and no \texttt{export()} call ever needs touching.

\subsection{Worked example 1: a US-state choropleth}\label{sec:gc-geo}
Evaluators hear one request more often than any other, \enquote{show me my states on a map,} and the geographic choropleth is how you answer it.\marginnote{A \emph{choropleth}\index{choropleth} shades each region by a value, so darker means more; the eye finds the hot spots without reading a single number.} We use the County Health Rankings \& Roadmaps (CHR) 2025 state extract, one row per state, and map the share of residents without health insurance. The geographic key is stored in \texttt{geo\_code} as a \enquote{US-XX} code (\texttt{US-TX}, \texttt{US-CA}), and \texttt{name()} points \texttt{googlechart} at it:
\par
{\footnotesize
\begin{verbatim}
import delimited using ///
    "`DATA'/chr_states_2025.csv", varnames(1) clear
googlechart uninsured, name(geo_code) type(geo)  ///
    georegion("US") georesolution("us-states")   ///
    scheme(blues)                     ///
    download datatable downloadpos(below)         ///
    title("Uninsured rate by state, 2025")        ///
    width(960) height(560) noopen                 ///
    export("`OUT'/01_geo_uninsured.html")
\end{verbatim}
}
Four options configure this chart type. The two geography settings, \texttt{georegion("US")} and \texttt{georesolution("us-states")}, tell the renderer to draw the fifty states plus DC rather than a world map. \texttt{scheme(blues)} sets the color ramp, so the darkest state is the highest uninsured rate. And \texttt{download datatable} adds two things a client asks for the moment they see the map: a menu to export the picture as an image, and a \enquote{view data} panel that shows and downloads the numbers behind it. The result is Figure~\ref{fig:gc-geo}: Texas is the darkest state at roughly nineteen percent uninsured, a cluster of southern and mountain-west states trails it, and the northeast is palest.

\begin{figure*}[htbp]
\centering
\includegraphics[width=\linewidth]{ch10_gc_geo.png}
\caption[A googlechart US-state choropleth]{The \texttt{type(geo)} choropleth of the uninsured rate by state, 2025, written by \texttt{ch10\_googlechart.do} with the command printed above. Darker is higher; the legend runs from 2.9 to 18.8 percent. In the browser this map is interactive: hover a state and its exact value appears. \texttt{googlechart}'s \texttt{geo} type stops at the state level; for a Texas county map, use the offline sibling \texttt{sparkta2} (Section~\ref{sec:sparkta2}).}
\label{fig:gc-geo}
\end{figure*}

Check the map against what you already know before you circulate it: the pattern is the familiar coverage gap between states that expanded Medicaid and states that did not, and Texas anchors the dark end of it every year, which is the external evidence that the pipeline drew the right map. The point for the suite is that this map is not a screenshot; it is written by the do-file from the CHR file, so when the 2026 release posts, rerunning the do-file redraws the map with no hand-editing. That is the replicable principle applied to an interactive graphic. The refresh is worth rehearsing before you need it: point the import at the prior year's file, rerun, and confirm the map redraws, because an update path that has never been exercised is an intention, not a fact.\marginnote{The \texttt{geo} type is the one most likely to blank behind a firewall, because it fetches boundary data at render time. Before you promise a client an air-gapped map, open the HTML with your network off and confirm it still draws; if it does not, ship the offline \texttt{sparkta2} choropleth or the static map of Chapter~\ref{ch:graphs}.}

\subsection{Worked example 2: an animated bubble with a Play button}\label{sec:gc-bubble}
When the story you are telling is a trajectory rather than a snapshot, motion can do real work. The animated bubble is \texttt{googlechart}'s Rosling-style chart, the moving-bubble format Hans Rosling's global-health talks made famous: each state is a circle, and pressing Play walks it through time so the reader watches the whole field move. We switch to the CHR \emph{panel} extract, which has one row per state per year from 2016 to 2025, because animation needs frames: without a row for each time value, there is nothing to animate.\marginnote{This is the one requirement newcomers miss. \texttt{time(year)} adds the Play button, but the button only does something if the data is a \emph{panel}, one row per entity per period. A single cross-section with \texttt{time()} draws a static bubble and no motion. Check that your data is long, not wide, before reaching for \texttt{time()}.} We plot children in poverty against premature death, size each bubble by how rural the state is, and color it by whether its uninsured rate clears ten percent:
\par
{\footnotesize
\begin{verbatim}
import delimited using ///
    "`DATA'/chr_states_panel.csv", varnames(1) clear
gen byte highunins = uninsured >= 10
label define hu 0 "Uninsured < 10%" 1 "Uninsured >= 10%"
label values highunins hu
googlechart child_poverty premature_death,       ///
    name(usps) over(highunins)                    ///
    sizevar(pct_rural) time(year) type(bubble)    ///
    tooltipvars(uninsured)                         ///
    download datatable downloadpos(below) ///
    title("Child poverty vs premature death"      ///
          " (press Play: 2016 to 2025)")           ///
    width(920) height(560) noopen                 ///
    export("`OUT'/04_bubble.html")
\end{verbatim}
}
Read the roles one at a time, because \texttt{type(bubble)} uses more of them than any other type. The two variables in the varlist are the x and y axes (poverty and premature death). \texttt{name(usps)} labels each bubble with its state postal code. \texttt{over(highunins)} colors the bubbles by group, so the high-uninsured states stand out. \texttt{sizevar(pct\_rural)} scales each circle by rurality. And \texttt{time(year)} is the option that adds the Play button and the year slider. The result is Figure~\ref{fig:gc-bubble}.

\begin{figure*}[htbp]
\centering
\includegraphics[width=\linewidth]{ch10_gc_bubble.png}
\caption[A googlechart animated bubble]{The \texttt{type(bubble)} chart with \texttt{time(year)}, written by \texttt{ch10\_googlechart.do}, shown at one frame. Each circle is a state, sized by how rural it is and colored by whether its uninsured rate clears ten percent. In the browser a Play button walks all fifty states from 2016 to 2025; here we print the final frame. Only 2024 and 2025 use observed CHR releases; earlier frames are simulated for the demonstration and do not document a historical trajectory.}
\label{fig:gc-bubble}
\end{figure*}

Use this animation to learn the controls, rather than to draw a conclusion about a decade of state trends. Only the 2024 and 2025 observations come from CHR releases; the earlier years were simulated to demonstrate movement through time. Keep that limitation in the visible chart note and any exported image. When adapting the example, replace the simulated rows with comparable observed releases and check whether changes in variable definitions or reference periods could create apparent movement. A static pair of observed years may be more informative when readers need a comparison rather than a demonstration of the Play control.

\subsection{Worked example 3: a searchable table}\label{sec:gc-table}
The other request an evaluator hears constantly is the opposite of a map: \enquote{let me find my own row.} The searchable table answers it. \texttt{type(table)} is the one type that usually takes \emph{no} varlist; instead you list the columns you want in \texttt{tooltipvars()}, and two switches turn an ordinary table into a tool: \texttt{tablesearch} adds a search box, and \texttt{tableheadersticky} freezes the header row so it stays visible as the reader scrolls:
\par
{\footnotesize
\begin{verbatim}
import delimited using ///
    "`DATA'/chr_states_2025.csv", varnames(1) clear
googlechart, type(table)                          ///
    tooltipvars(state uninsured median_income     ///
                premature_death obesity           ///
                child_poverty unemployment        ///
                pct_rural)                         ///
    tablesearch tableheadersticky                 ///
    download datatable downloadpos(below) ///
    title("2025 County Health Rankings:"          ///
          " state measures (searchable)")          ///
    width(980) height(460) noopen                 ///
    export("`OUT'/05_table.html")
\end{verbatim}
}
The payoff is in Figure~\ref{fig:gc-table}: a client opens the page, types their state name into the search box, and reads their own row without a spreadsheet, or clicks a column header to sort every state by that measure. That is the accessible principle served directly, and because the file is written by the do-file it refreshes when the CHR data does.\marginnote{Column choice is the planning decision: lead \texttt{tooltipvars()} with the identifier the reader will search on and the two or three measures they came for. If a table opens on eight columns of equal weight, the reader has to do the triage the builder skipped.}

\begin{figure*}[htbp]
\centering
\includegraphics[width=\linewidth]{ch10_gc_table.png}
\caption[A googlechart searchable table]{The \texttt{type(table)} output with \texttt{tablesearch} and \texttt{tableheadersticky}, written by \texttt{ch10\_googlechart.do}. All 51 rows (fifty states plus DC) and eight measures; the search box at top filters as the reader types, the sticky header keeps the column names in view, and any column sorts on a click. The columns come from \texttt{tooltipvars()}, not a varlist, which is the one syntax quirk of the table type.}
\label{fig:gc-table}
\end{figure*}

\subsection{Worked example 4: a diverging stacked bar}\label{sec:gc-divbar}
The last type is the one survey work needs most and general tools handle worst: the diverging stacked bar\index{diverging stacked bar} for Likert data. A \emph{Likert}\index{Likert scale} item is the familiar strongly-disagree-to-strongly-agree scale, and the honest way to chart it is to center the neutral response on zero, push disagreement left and agreement right, so the balance of opinion reads at a glance. \texttt{type(divbar)} does exactly this. It needs two roles: \texttt{name()} for the question and \texttt{level()} for the response category, plus \texttt{levelorder()} to fix the stack order and \texttt{centerlevel()} to name the neutral level that sits on zero:
\par
{\footnotesize
\begin{verbatim}
googlechart share, name(q) level(response)        ///
    type(divbar)                                  ///
    levelorder("Strongly disagree|Disagree"       ///
        + "|Neutral|Agree|Strongly agree")         ///
    centerlevel(Neutral)                           ///
    download datatable downloadpos(below) ///
    title("Texans on state policy"                ///
          " (illustrative survey data)")           ///
    subtitle("Diverging stacked bar;"             ///
             " neutral centered on zero")          ///
    width(1040) height(420) noopen                ///
    export("`OUT'/06_divbar.html")
\end{verbatim}
}
The data here is a small synthetic set, three policy statements each with five Likert levels, written inline in the do-file so the shape is illustrative rather than a real poll. Figure~\ref{fig:gc-divbar} shows the result: each statement is one horizontal bar split at zero, disagreement in reds to the left and agreement in blues to the right, with the neutral slice straddling the center line.

\begin{figure*}[htbp]
\centering
\includegraphics[width=\linewidth]{ch10_gc_divbar.png}
\caption[A googlechart diverging stacked bar]{The \texttt{type(divbar)} chart, written by \texttt{ch10\_googlechart.do} from illustrative Likert data. \texttt{centerlevel(Neutral)} puts the neutral response on the zero line and sign-flips the disagree side to the left, so a reader can compare the balance of opinion item by item. This is the Pew-style diverging bar; the underlying numbers are synthetic and shown only for shape.}
\label{fig:gc-divbar}
\end{figure*}

Between them, the geo map, the animated bubble, the searchable table, and the diverging bar cover most of what a board actually asks for, and one do-file writes all four offline with no key. Two options recur across all four. \texttt{scheme()} sets the color ramp in one word\marginnote{In \texttt{scheme()}: \texttt{scheme(blues)} for a sequential map, \texttt{scheme(rdbu)} for a diverging one.}, so every page from your shop can use the same palette without per-chart color fiddling. And \texttt{download datatable} pairs an image-export menu with a \enquote{view the data} panel, so the client takes away both the picture and the numbers. The other chart types follow the same grammar: a \texttt{type()}, a varlist or a set of role options, and an \texttt{export()}.

With fourteen types on the menu, the temptation is to build several and let the client choose. You will get there faster from the other direction: name the question first, then pick the type that answers it. \enquote{Where are we} is the geo map, \enquote{how did we get here} is the line or the bubble, \enquote{where is my row} is the searchable table, and three pages that each answer a named question beat ten that answer none in particular. The types you did not build can wait for the day a stakeholder asks a question shaped like them.

\section{Static, interactive, or animated: prefer the simplest form that delivers the finding}
Does this deliverable need to move at all? Ask before you build, because choosing interactivity is making a claim about how the audience will use the evidence, and the claim should be chosen, not defaulted to. In the two subsections that follow we set out the test we run first: place the deliverable on the author-driven to reader-driven spectrum, then read how each form fails, so you can say in one sentence why this page is static rather than interactive when a client asks why the map does not move. A reader skimming a printed report or a slide wants a clean static figure; a client who will genuinely explore, filter to their county, compare their site, wants interactivity; a change over time that the eye should follow wants animation, but only when the motion adds information rather than decoration. Animation that adds nothing subtracts, because it makes the reader wait for a story a static small-multiple would have told at once.\sidenote{Rosling's bubble charts are the honorable exception: motion works there because the story \emph{is} the trajectory over time. If your animation would read just as well as a set of before/after panels, ship the panels.}

\subsection{Which narrative structure for a board versus an explorer}\label{ch10-sec:structures}
Structure comes before form: decide whether you or the reader steers, and by the end of this subsection you will be able to name the structure you are building and the venue it belongs in before the first line of chart code. \textcite{segel2010narrative} studied dozens of narrative graphics and asked, of each one, who holds the controls. Their answer is a spectrum, author-driven at one end (the maker fixes the message), reader-driven at the other (the reader takes the controls and finds their own path), and every interactive feature you add slides the deliverable toward the reader's end. Their middle ground supplies this chapter's working vocabulary: the \emph{martini-glass}, which walks the reader down the author's path first and then opens into free exploration, and the \emph{drill-down story}, which inverts the order, opening on a general view and letting the reader pick the case to expand.\marginnote{The third \textcite{segel2010narrative} structure, the \emph{interactive slideshow}, keeps the author's ordering across steps while allowing interaction inside each one; it fits a walked-through talk more than a left-alone web page.} Figure~\ref{fig:ch10-viz-spectrum} places the author-driven slide and both structures on the spectrum, each paired with the venue it fits, and turns the design choice into the utilization-focused question: who will act on this finding, and how much of the drill-down should they own?

\begin{figure*}[htbp]
\centering
\resizebox{\linewidth}{!}{%
\begin{tikzpicture}[node distance=0.6cm and 1.5cm]
  \node[stage] (auth) {Author-\\driven};
  \node[stage, right=of auth] (mart) {Martini-\\glass};
  \node[stage, right=of mart] (drill) {Drill-\\down};
  \draw[flow] (auth) -- (mart);
  \draw[flow] (mart) -- (drill);
  \node[annot, left=0.15cm of auth, text width=1.6cm, align=right]
    {maker fixes\\the message};
  \node[annot, right=0.15cm of drill, text width=1.6cm, align=left]
    {reader picks\\the path};
  \node[substep, below=0.5cm of auth, text width=2.0cm]
    {One path, no controls\\\emph{Full-board slide}};
  \node[substep, below=0.5cm of mart, text width=2.0cm]
    {Headline first,\\then roam\\\emph{Briefing}};
  \node[substep, below=0.5cm of drill, text width=2.0cm]
    {Open on all,\\expand one\\\emph{Manager portal}};
\end{tikzpicture}%
}
\caption[Narrative structures on the control spectrum]{The three narrative structures of \textcite{segel2010narrative} arranged left to right as control passes from maker to reader, each tagged with the venue it fits. Read it as a ladder: hand a full board the author-driven slide, a briefing the martini-glass that states the headline before it opens, and a manager portal the drill-down that opens on the whole field.}
\label{fig:ch10-viz-spectrum}
\end{figure*}

The judgment, then, is not \enquote{interactive or not} but \enquote{which structure for which reader.} A board that meets for ninety minutes and needs one decision is an author-driven or martini-glass audience: give them the message first and let exploration follow the point, never precede it. An analyst who will live in the data for a week is a drill-down audience: open on the whole field and let them expand their own county without your narration in the way. The mistake an embedded evaluator makes under deadline is to ship a drill-down structure to a board, which buries the finding behind a click the board will not make, or an author-driven figure to an explorer, which locks the analyst out of the question they came to ask.\marginnote{The searchable table of Section~\ref{sec:gc-table} is a drill-down surface; the spark wall of Section~\ref{sec:sparkwall} paired with one interactive map is a martini-glass deliverable, headline first, exploration second.} Naming the structure before you build tells you which one you are actually making.

The literature attaches three cautions to that delegation, and each compresses to a sentence. Consider the defaults first. \textcite{hullman2011rhetoric} show that an interactive's opening view, ordering, and annotations steer interpretation the way framing steers a survey item, so make the defaults defensible and let the reader escape them, which the download-and-datatable menu of Section~\ref{sec:gc-table} provides. Delegation is a hypothesis, not a delivery: in field experiments on live news and gallery sites, \textcite{boy2015storytelling} found that narrative and interactive hooks did not reliably increase how much readers actually explored, so build the drill-down only where you have reason to think the stakeholder will drill. And every control costs a click: \textcite{few2013dashboard}'s rule that a monitoring surface must read at a glance means a slider that hides the headline behind an action a busy manager will not take subtracts rather than adds. Treat presentation as a first-class part of the analysis rather than an afterthought \parencite{kosara2013storytelling}, and the three cautions become design inputs instead of reviewer objections.

That reasoning turns into a build step you can finish at a desk before the first \texttt{googlechart} call. Enumerate the drill-downs your users will actually perform, and there are rarely more than three: filter to my county, trace the measure over time, compare against the target. Write each one down with a name attached, \enquote{the Travis County program manager filters to Travis at the March meeting,} because a drill-down without a named clicker is a hypothesis with no one to test it. Build only the paths on the list; every control beyond them is surface area that can confuse a reader and break at refresh time. And when the list comes back empty, when no one on the distribution can be named as the person who will click, the exercise has answered the static-versus-interactive question for you: ship the static panel and keep the interactive budget for a deliverable that has a clicker.

\subsection{Match the form to the venue, and read the failure column first}\label{ch10-sec:forms}
We suggest you match the medium to the audience and the venue, and prefer the simplest form that communicates the finding. Table~\ref{tab:interactive-choice} lays out the three forms against the four questions that actually decide between them: who the audience is, where the deliverable lives, how it updates, and the way each form fails. Reading the failure column is the most useful thing you can do before committing, because every form fails differently and the wrong failure is expensive.

\begin{table*}[htbp]
\centering
\caption[Choosing static, interactive, or animated]{A decision table for the three delivery forms. Match the row to your audience and venue, then read the failure column before you commit: each form breaks in a different way, and the cheapest mistake to avoid is shipping motion or interactivity that the venue cannot render.}
\label{tab:interactive-choice}
{\footnotesize
\begin{tabular}{@{}llllp{6.0cm}@{}}
\hline
Form & Audience & Venue & Update & Failure mode \\
\hline
Static & Skimmer & Print, PDF, slide & Rerun do-file & Unreadable, misleading, or stale chart \\
Interactive & Explorer & Web page, portal & Rebuild HTML & Blanks behind a firewall \\
Animated & Watcher & Screen, talk & Re-render frames & Motion adds no signal \\
\hline
\end{tabular}
}
\end{table*}

Read the table as a ladder of risk. Static output avoids a live drawing-library dependency, but it still needs readable labels, an appropriate scale, and a visible data date. When you choose interactive output, you buy exploration at the cost of a live rendering dependency, so the page fails exactly where Section~\ref{sec:sparkwall}'s margin notes warn, behind a firewall or when a CDN is blocked. Animated output buys a followed-over-time story at the cost of the reader's patience, so it fails when the motion adds no information the eye needed to see move. Choose deliberately and the interactive output stays in service of the message; default to it and you end up demonstrating the tool rather than delivering the finding, which is the temptation the accessible principle warns against.

A practical corollary follows from the failure column. When you are not sure the venue will render an interactive or animated deliverable, build the static version first and treat the fancier form as an enhancement, not the base case.\marginnote{This is progressive enhancement, borrowed from web design: ship a version that works everywhere, then layer on features for the environments that support them. A dashboard that degrades to a readable table has a floor; one that degrades to a blank page has none.} The spark wall of Section~\ref{sec:sparkwall} is precisely this hedge: a static small-multiple that answers the board's question on its own, so if the interactive map blanks behind the client's firewall, the exhibit that survives still answers it. And when the venue is offline for good, an air-gapped agency or an embargoed report, the choice inside the interactive row is \texttt{sparkta2} over \texttt{googlechart}, because only the offline sibling stores its drawing engine in the file. Keep that ordering, static as the floor and interactivity as the bonus, and a blocked CDN costs you a feature rather than the deliverable.

\section{Build the chart for the readers you actually have}\label{ch10-sec:a11y}
Look around the boardroom before you pick a palette.\marginnote{Roughly one man in twelve and one woman in two hundred has a red-green color vision deficiency, so a board of thirty likely seats one or two members who read no signal at all from a red-versus-green diverging bar.} The suite's accessible principle asks that an audience read the finding without a specialist to translate it, and for interactive output that is not a metaphor: a share of your real readers cannot see the chart the way you drew it, so accessibility is a design constraint you either meet or fail. Perception researchers treat color as a channel with limits, not a free variable \parencite{franconeri2021science}, and the accessibility field has turned those limits into rules you can follow before shipping. We work through three of them here, a colorblind-safe palette, a contrast floor you can cite rather than eyeball, and the numbers delivered beside the picture, and fold all three at the end into a ten-minute pass you run on every page you build.

\subsection{Color, contrast, and the numbers beside the picture}\label{ch10-sec:a11y-rules}
The first rule is about color: it is worth picking a palette that survives color vision deficiency, because the field has already settled on an answer you can adopt. \textcite{okabeito2008cud} designed the Color Universal Design palette, eight qualitative colors chosen so protanopes, deuteranopes, and tritanopes, people with the three main forms of color-vision deficiency, can still tell them apart, and \textcite{wong2011color} brought it into scientific practice as the Nature Methods default. Table~\ref{tab:ch10-viz-cud} prints the eight colors with the hex codes you paste straight into a \texttt{scheme()} or a style file.

\begin{table*}[htbp]
\centering
\caption[The Okabe--Ito colorblind-safe palette]{The eight Color Universal Design colors of \textcite{okabeito2008cud}, each with the hex code to hand a chart. Chosen so the common forms of color vision deficiency can still separate them; use these in place of an arbitrary categorical ramp, and let position or luminance, not hue alone, show the finding. Viewed in grayscale the swatches still differ in lightness, which is the redundant-channel rule made visible.}
\label{tab:ch10-viz-cud}
{\definecolor{cudBlack}{HTML}{000000}%
\definecolor{cudOrange}{HTML}{E69F00}%
\definecolor{cudSky}{HTML}{56B4E9}%
\definecolor{cudGreen}{HTML}{009E73}%
\definecolor{cudYellow}{HTML}{F0E442}%
\definecolor{cudBlue}{HTML}{0072B2}%
\definecolor{cudVermillion}{HTML}{D55E00}%
\definecolor{cudPurple}{HTML}{CC79A7}%
\newcommand{\cudswatch}[1]{\fcolorbox{black!30}{#1}%
  {\phantom{\rule{1.6em}{0.9ex}}}}%
\footnotesize
\begin{tabular}{lll}
\hline
Swatch & Name & Hex \\
\hline
\cudswatch{cudBlack}      & Black          & \texttt{000000} \\
\cudswatch{cudOrange}     & Orange         & \texttt{E69F00} \\
\cudswatch{cudSky}        & Sky blue       & \texttt{56B4E9} \\
\cudswatch{cudGreen}      & Bluish green   & \texttt{009E73} \\
\cudswatch{cudYellow}     & Yellow         & \texttt{F0E442} \\
\cudswatch{cudBlue}       & Blue           & \texttt{0072B2} \\
\cudswatch{cudVermillion} & Vermillion     & \texttt{D55E00} \\
\cudswatch{cudPurple}     & Reddish purple & \texttt{CC79A7} \\
\hline
\end{tabular}}
\end{table*} On a \texttt{googlechart} or \texttt{sparkta2} page, we suggest you feed \texttt{scheme()} a colorblind-safe ramp rather than an arbitrary one, and prefer a sequential or diverging scheme, a single hue darkening or two hues meeting at a neutral middle, over a rainbow, because a luminance gradient reads even when hue does not. The spark wall of Section~\ref{sec:sparkwall} is already safe on this count, because it encodes its story in line \emph{position} rather than color; a reader who sees no color still reads six rising trajectories, which is the deeper rule that color should be a redundant channel, never the only one.\marginnote{The test costs nothing: view your finished page through a color-blindness simulator, or desaturate a screenshot to grayscale, and check whether the finding still reads. If the categories collapse into the same gray, hue was doing work that shape, position, or luminance should have shared.}

The second rule is sufficient contrast, and here the web has a published standard the evaluator can cite rather than eyeball. The Web Content Accessibility Guidelines set a floor of a 4.5:1 luminance-contrast ratio for normal text and 3:1 for the graphical objects a reader must see to understand the chart \parencite{w3c2018wcag}, a threshold a free contrast checker measures for you against the two colors that meet. The ratio itself is a simple object:
\[
  C \;=\; \frac{L_{\text{light}} + 0.05}{L_{\text{dark}} + 0.05},
  \qquad 1 \le C \le 21,
\]
where $L$ is each color's relative luminance, a weighted sum of its red, green, and blue channels that runs from $0$ for black to $1$ for white, and the lighter color goes on top. The rule is then just $C \ge 4.5$ for text and $C \ge 3$ for chart furniture\marginnote{The constant $0.05$ keeps the denominator positive even for black; the ratio maxes at $21{:}1$ for pure black on pure white.}. For an interactive page this bites in the furniture as much as the data: a pale-gray axis label on white, a light tooltip on a light map, or a legend swatch that fails 3:1 against its background each quietly locks out the low-vision reader the accessible principle means to include. Checking normal-size text labels against 4.5:1 and essential graphical marks against 3:1 is a five-minute pass that keeps a page's promise to the reader who needs it most.

The third rule is that a chart is more than its pixels, so an interactive page should not strand a reader who navigates by keyboard or hears the page through a screen reader. Two habits do most of the work. Ship the numbers alongside the picture, which is what the \texttt{download datatable} panel of Section~\ref{sec:gc-table} already does: a screen-reader user who cannot perceive the map still reads and sorts the underlying table, so the datatable is an accessibility feature, not only a download convenience. And give the page a text alternative that states the finding in a sentence, so a reader who never sees the chart still gets its message. These are not exotic requirements; they ask the graphic to \enquote{say the finding in words}, because words reach the readers that pixels cannot. If part of the audience cannot see, click, or hear an interactive deliverable, it is not accessible in the suite's sense, however polished it looks.

Consider folding these checks into a single pre-ship pass rather than treating each as a separate virtue. Ten minutes covers the chapter's whole exposure: open the page with the network off (does it draw?), view a screenshot in grayscale (does the finding survive?), check normal-size text against 4.5:1 and essential graphical marks against 3:1, and confirm the datatable accompanies the picture. Run the same four checks on every page, in the same order, and accessibility stops being an aspiration and becomes a routine the shop does not have to remember to have.

\section{A live dashboard is a promise to keep the pipeline alive}\label{ch10-sec:maintenance}
Interactivity creates a cost the static figure does not, and the cost comes due after you ship rather than before. A printed spark wall is finished the moment the do-file runs; a live interactive page, by contrast, sets an expectation that its data is current, so the day you hand a board a dashboard you have quietly promised to keep feeding it. That promise is why we attach one caveat to every interactive deliverable in this chapter: the graphic is only as trustworthy as the pipeline behind it, and a dashboard that silently shows last year's numbers is worse than a dated PDF, because the PDF shows its date on its face while the dashboard looks live and misleads.

Three failure modes make the promise concrete. The first is data staleness. The QCEW loop of Section~\ref{sec:sparkwall} pulls from live URLs, so a page built once and left alone drifts out of date the moment a new quarter posts, and a reader who trusts the freshness the interactivity implies is misled. You already know the replicable remedy from the rest of this chapter: keep the deliverable a do-file rather than a hand-made artifact, so a refresh is one rerun rather than a rebuild, and schedule the rerun rather than waiting for someone to notice the numbers went old. The second is dependency rot. An online page like \texttt{googlechart}'s links out to a rendering library it does not contain (Section~\ref{sec:gc-how}), and a library that changes its interface, moves its URL, or is withdrawn can blank a page that worked for years. Vendoring the engine inside the file, which is what the offline sibling \texttt{sparkta2} does, trades a live dependency for a frozen one, which removes that external dependency; browser changes and defects in the bundled code can still affect the page. The third is orphaned ownership. A dashboard built by an evaluator who then rolls off the contract becomes a liability the agency cannot repair, a broken page with the agency's name on it and no one left who knows the code.

Those three failure modes are sobering rather than discouraging: match the deliverable's lifespan to the maintenance you can actually promise. A one-time interactive for a single meeting needs no upkeep and creates no risk, so build it freely. A standing portal a board will consult monthly for two years is a maintenance commitment, so either budget the reruns and the version pins that keep it alive, or ship the offline, frozen \texttt{sparkta2} version that degrades gracefully rather than the online one that decays silently.\marginnote{\textcite{few2013dashboard} frames a monitoring surface as something a reader returns to, which is why its data must stay current to keep faith with that return; a dashboard read once and never refreshed has broken its at-a-glance promise.} The extensible principle cuts both ways here. The same do-file pipeline that lets you add a county also lets a successor refresh the page after you leave, but only if you shipped the code and the documentation, not just the rendered HTML. Hand over the do-file, the data source, and a one-line \enquote{rerun this each quarter} note, and the dashboard stops being a promise only you can keep; the agency can now keep it without you.

That handover works best when you settle the refresh calendar at ship time, while everyone who built the page is still in the room. Three decisions make it real. Who reruns: a named person, not a role, because \enquote{the data team} is how a dashboard orphans itself. On what trigger: a calendar date or a data-release event, and for the QCEW pipeline of Section~\ref{sec:sparkwall} the natural trigger is the quarterly posting, which lags the quarter by several months, so the note names the expected posting month rather than the quarter's end. What gets archived: each superseded HTML saved with a date suffix before the new build overwrites it, so \enquote{what did the board see in March} still has an answer when a number changes underneath a decision. Written on one page and handed over with the do-file, the calendar is the difference between a dashboard that decays and one that merely waits for its next rerun.

When the trigger is a calendar date, consider taking the rerun out of human memory entirely. The same scheduler that pulled survey responses every morning in Chapter~\ref{ch:surveys} can run a chart build: one \texttt{cron} line on macOS or Linux, or one Task Scheduler job on Windows, runs the do-file in batch on the expected posting date, writes the rebuilt pages, and leaves a dated log. The named person's job then shrinks from remembering the rerun to reading the log and clicking through the refreshed pages once, ten minutes on a known morning instead of a task that competes with everything else for attention. Keep the human in that last step rather than automating it away: the scheduler can rebuild a page, but only a person notices that a county's wage series suddenly halved because the source revised its history, and a dashboard that republishes bad data on schedule is automation working against you.

\begin{appliedexample}[One dataset, three deliverables]\label{ae:ch10-threedeliv}
\textbf{Scenario.} The workforce board wants the county wage story in three places: a one-page printed brief, an exploratory portal on the agency website, and a two-minute clip in the director's talk.

\medskip
\textbf{The build.} For the county brief, use the static spark wall built by \texttt{ch10\_smallmult.do}. An interactive table could present the same county-quarter values, provided you adapt the table example's columns and labels to those data. The \texttt{googlechart} geographic example in this chapter uses a different, state-level CHR file and cannot produce a county map from the QCEW panel. For a county map, use an available mapping tool with county boundaries or deliver a static map; \texttt{sparkta2} remains a proposed option until its source is released. Keep the state animation as a separate teaching demonstration because its earlier frames are simulated. Document the source and build command for each deliverable so that a colleague can refresh the intended dataset.
\end{appliedexample}

\paragraph{Where this leaves you.} This chapter opened with a board asking for \enquote{something like the newspaper does} and closes with you able to build it without leaving the do-file: spark walls for scanning a portfolio, interactive maps and searchable tables, an animated bubble for the one story that really is a trajectory, each written by code that reruns when the data refreshes. The judgments you picked up along the way matter more than any single command. Before the first \texttt{googlechart} call, you place the deliverable on the author-driven to reader-driven spectrum and write down the drill-downs a named stakeholder will actually perform, treating interactivity as a hypothesis about a clicker rather than a default. You build the static floor first, so a blocked CDN costs a feature rather than the deliverable. You check the palette, the contrast, and the shipped numbers so the readers in the actual room can read the page. And at ship time you write the refresh calendar, who reruns, on what trigger, what gets archived, then hand the rerun itself to a scheduler and keep only the looking, because a live dashboard is a promise to keep its pipeline current and a promise is easier to keep when a machine does the remembering. Chapter~\ref{ch:sheets} turns to the format most clients actually live in, the spreadsheet, and how to deliver through it without giving up the pipeline.

\section*{Exercises}\addcontentsline{toc}{section}{Exercises}
\begin{enumerate}
\item \emph{Try it on your own data.} Take a dataset you already report on that has one outcome measured over time for several sites, build a spark wall with \texttt{twoway line} and \texttt{by()}, strip the axes as Section~\ref{sec:sparkwall} does, and confirm every site's panel actually drew by checking the panel count against your number of sites.
\item Rerun \texttt{ch10\_smallmult.do} after adding two more Texas counties to the \texttt{counties} local, then report whether the two new panels share the same first-quarter bonus spike the original six show.
\item Predict, before you run it, how many of the 144 QCEW downloads in Section~\ref{sec:sparkwall} will land for your county set; rerun the loop, check \texttt{\_N} against your prediction, and explain any gap by the newest quarters that have not yet posted.
\item Break the FIPS loop on purpose: drop the \texttt{capture} and \texttt{if \_rc continue} guard, rerun the download over a range that includes an unposted quarter, and confirm the loop now aborts on the first missing file instead of skipping it.
\item Rerun \texttt{ch10\_googlechart.do} to rebuild the \texttt{type(geo)} choropleth, open the resulting HTML with your network turned off, and confirm the map blanks; then swap \texttt{scheme(blues)} for \texttt{scheme(rdbu)} and describe how the color ramp changes which states read as extreme.
\item \emph{Place a deliverable on the control spectrum.} Take one interactive page you built above and locate it on the author-driven to reader-driven spectrum of \textcite{segel2010narrative}: name the default view it opens on, name the drill-down it delegates to the reader, and name the stakeholder who would actually perform that drill-down. If you cannot name that stakeholder, rebuild the finding as a static panel and say what the interactivity was costing without buying.
\item \emph{Check one page for accessibility.} Take an interactive page you built above, desaturate a screenshot of it to grayscale (or run it through a color-blindness simulator), and say whether the finding still reads once hue is gone; then check normal-size axis and legend text against 4.5:1 and essential graphical marks against 3:1 of Section~\ref{ch10-sec:a11y}, and name one change, a colorblind-safe \texttt{scheme()}, a darker label, or the \texttt{download datatable} panel, that would let a channel-limited reader use the page.
\item \emph{Write the refresh calendar.} For the portal deliverable in the applied example above, draft the one-page refresh note of Section~\ref{ch10-sec:maintenance}: name the person who reruns \texttt{ch10\_smallmult.do}, the trigger that prompts the rerun, and the file-name pattern that archives the superseded HTML; then confirm the next-quarter refresh is a one-line change by finding the single header line the rerun requires you to edit.
\end{enumerate}
